% !TEX root = ../main.tex
% !TEX program = lualatex
\documentclass[../main.tex]{subfiles}
\IfSubfilesClassLoaded{\externaldocument{\subfix{../build/main}}}{}
% ====================
% File: 14_EVM.tex
% ====================
\begin{document}
\IfSubfilesClassLoaded{\chapter{EDO Study Guide}}{}

% === EVM (EDO 3.6.1 & 3.6.2) ===
\section{Earned Value Management (EVM)}

\subsection{Summary}
\begin{description}
  \item[\textbf{Purpose.}] \ac{evm} integrates scope, schedule, cost, technical progress, and risk against the contract
  execution baseline so the \ac{pmo} can see performance trends early.
  \item[\textbf{Baseline cue.}] The \ac{pmb} is the contractor's time-phased execution baseline for authorized work; it is
  not the \ac{apb}, which is the government program-level cost, schedule, and performance agreement with the \ac{mdauthority}.
  \item[\textbf{Board math.}] \ac{pv}/\ac{bcws}, \ac{ev}/\ac{bcwp}, and \ac{ac}/\ac{acwp} drive variance, index,
  forecast, and \ac{tcpi} answers; always state whether the issue is cost, schedule, or baseline realism.
\end{description}
~\srcCite{edo-3-6-1-evm-intro-2025,edo-3-6-2-evm-analysis-2025,DoD-IPM-EVM-Definitions-2026}

\subsection{What EVM is}

\ac{evm} integrates technical scope, schedule, resources/cost, and risk in a baseline plan against which progress is
measured. In \ac{dod} use, it supports management decisions by showing cost, schedule, and technical performance trends
against the contract execution baseline~\autocite{DoD-IPM-EVM-Definitions-2026}.

\subsection{Key Components of EVM}
Core cost and schedule metrics measure performance against the \ac{pmb}; they do not establish the baseline by
themselves~\autocite{edo-3-6-1-evm-intro-2025,edo-3-6-2-evm-analysis-2025,DoD-IPM-EVM-Definitions-2026}.
\begin{description}
  \item[\textbf{PV / BCWS.}] The budgeted cost for the work scheduled to be completed by a specific date.
  \item[\textbf{EV / BCWP.}] The budgeted cost for the work actually completed by a specific date.
  \item[\textbf{AC / ACWP.}] The actual cost incurred for the work completed by a specific date.
  \item[\textbf{BAC.}] The total budget allocated for the project.
  \item[\textbf{EAC.}] The forecasted total cost of the project based on current performance.
  \item[\textbf{CPI.}] A measure of cost efficiency, calculated as $\text{CPI} = \text{EV} / \text{AC}$.
        A \ac{cpi} less than 1 indicates a cost overrun.
  \item[\textbf{SPI.}] A measure of schedule efficiency, calculated as $\text{SPI} = \text{EV} / \text{PV}$.
        A \ac{spi} less than 1 indicates a schedule delay.
  \item[\textbf{SV.}] The difference between earned value and planned value, calculated as
        $\text{SV} = \text{EV} - \text{PV}$. A negative \ac{sv} indicates a schedule delay.
  \item[\textbf{CV.}] The difference between earned value and actual cost, calculated as
        $\text{CV} = \text{EV} - \text{AC}$. A negative \ac{cv} indicates a cost overrun.
\end{description}
\Info{\ac{pv}, \ac{ev}, and \ac{ac} are the modern labels; \ac{bcws}, \ac{bcwp}, and \ac{acwp} are the equivalent legacy EVM terms.}

\subsection{When to Use EVM}
\ac{evm} is mandated on cost or incentive contracts (and applicable subcontracts) when the program meets the policy thresholds and the effort is discretely measurable~\autocite{edo-3-6-1-evm-intro-2025,DFARS-EVMS-Deviation-2026}. Key decision points for the \ac{pm} include:
\begin{description}
  \item[\$100 million and above.] Require an \ac{evms} that the cognizant Federal agency has determined complies with
        \ac{eia}-748.
  \item[\$50 million to \$99.9 million.] Require an \ac{evms} for cost or incentive contracts and subcontracts, but a
        formal compliance determination is not required unless separately directed.
  \item[Below \$50 million.] Apply \ac{evm} only when the \ac{pm} determines it is useful; below \$20 million, the
        cost-benefit analysis must be documented in the contract file.
\end{description}
Additional policy considerations follow~\autocite{edo-3-6-1-evm-intro-2025,DoDI5000-85,DFARS-EVMS-Deviation-2026}.
\begin{itemize}
  \item Applicability determinations must confirm the work scope can be discretely measured before mandating \ac{evm}.
  \item \ac{evm} on \ac{ffp}, time-and-materials, or \ac{loe} contracts is discouraged absent clear management benefit.
  \item Contracts primarily for software, including \ac{dodi}~5000.87 software-acquisition pathway contracts, are exempt from \ac{dfars} Subpart~234.2 under the active class deviation.
  \item Contracts that require \ac{evm} normally also require \ac{ipmdar} data and an \ac{ibr} as early as practicable
        and no later than 180 calendar days after contract award, exercise of a significant option, or major modification.
\end{itemize}
\subsection{EVM Compliance}
\ac{evm} clauses are flowed in the solicitation and award package. Under Class Deviation~2026-O0011, use \ac{dfars}~252.234-7998 and 252.234-7999 instead of the legacy 252.234-7001 and 252.234-7002 text~\autocite{DFARS-EVMS-Deviation-2026}. Compliance expectations follow~\autocite{edo-3-6-1-evm-intro-2025,DoDI5000-85,DFARS-EVMS-Deviation-2026}.
\begin{itemize}
  \item The program office integrates \ac{evm} planning into the \ac{wbs}, maintains the \ac{pmb}, and ensures timely receipt and analysis of \ac{ipmdar} submissions.
  \item When \ac{dod} is the cognizant Federal agency, \ac{dcma} determines \ac{evms} compliance against \ac{ansi}/\ac{eia}-748; the contracting officer approves or disapproves the system with input from functional specialists and auditors.
  \item Surveillance access is contractual: the applicable \ac{dfars} clause gives the Government access to contractor processes, data, and personnel needed to assess ongoing \ac{evms} compliance.
  \item Department and service focal points such as \ac{dod} Integrated Program Management within Acquisition Policy
        and Innovation and the \ac{don} \ac{cevm} support applicability determinations, tailoring,
        policy interpretation, and waiver/deviation coordination.
  \item \ac{supship} can execute validation and surveillance responsibilities for shipbuilding programs on behalf of \ac{dcma}\FlagBilletNote{Current \ac{supship} and \ac{navsea} reserve flag billet context}.
\end{itemize}

\subsection{EVM Principal Players}
Programmatic, contracting, and oversight roles synchronize to sustain disciplined performance
management~\autocite{edo-3-6-1-evm-intro-2025,DoD-IPM-EVM-Definitions-2026,DoD-IPM-Policy-Guidance-2026,DON-CEVM-2026}.
\begin{description}
  \item[\textbf{PMO}:] Implements \ac{evm} on the contract, ensures solicitations and awards contain the required clauses, and works with the contracting activity to tailor reporting and \ac{ims} requirements.
  \item[\textbf{DCMA}:] When \ac{dod} is the cognizant Federal agency, determines contractor \ac{evms} compliance, supports acceptance and surveillance, maintains the official acceptance roster, and analyzes \ac{ipmdar} and schedule submissions to focus government attention on emerging issues.
  \item[\textbf{DCAA}:] Conducts accounting system audits, corroborates cost data, and supports \ac{dcma} and the contracting officer during surveillance events and reviews.
  \item[\textbf{DoD Integrated Program Management}:] Serves as the department focal point for \ac{evm} policy, guidance, reporting, and competency management through Acquisition Policy and Innovation.
  \item[\textbf{Naval CEVM}:] Provides \ac{don} policy, guidance, and tools for applying \ac{evm}; coordinates applicability and reporting questions with \ac{dasnapb}, \ac{osd}, and program offices.
  \item[\textbf{SUPSHIP}:] Performs many \ac{dcma}/\ac{dcaa} surveillance roles for shipbuilding contracts when delegated by the cognizant program office.\FlagBilletNote{Current \ac{supship} flag billet holders.}
\end{description}

\subsection{Performance Measurement Baseline (PMB)}

\Fig{\ref{fig:pmb}} shows the characteristic S-curve depicting cumulative \ac{pv}/\ac{bcws}, while
\Fig{\ref{fig:pmb_dev}} summarizes the progression from defining the \ac{sow} to final \ac{pmb} adjustments. \ac{pmb}s are:
\begin{itemize}
  \item Scoped, scheduled, and budgeted plans for the authorized work
  \item Time-phased budgets that align to the master schedule
  \item The basis for cost and schedule performance management
  \item The source of planned value; cumulative \ac{pv} is the \ac{pmb} budget scheduled through a status date
\end{itemize}
Do not confuse the \ac{pmb} with the Acquisition Program Baseline. The Acquisition Program Baseline is the
program-level cost, schedule, and performance agreement with the Milestone Decision Authority; the \ac{pmb} is the
contract execution baseline used to measure earned value performance.
\Info{Board trap: the \ac{ibr} assesses \ac{pmb} realism, risk, schedule logic, resource adequacy, and complete scope coverage. It is not the same as \ac{mdauthority} approval of the \ac{apb}.}

\begin{figure}[H]
  \centering
  \includegraphics[width=0.5\linewidth]{PMB.png}
  \caption[PMB ``S''-Curve]{\ac{pmb} cost vs.\ time chart showing the ``S''-curve. \srcCite{edo-3-6-1-evm-intro-2025}}
  \label{fig:pmb}
\end{figure}

\begin{figure}[H]
  \centering
  \includegraphics[width=0.85\linewidth]{pmb_development.png}
  \caption[PMB Development Flow]{The development flow for \ac{pmb}. \srcCite{edo-3-6-1-evm-intro-2025}}
  \label{fig:pmb_dev}
\end{figure}

\subsubsection{PMB Development}
\begin{enumerate}
  \item[Step 1:] \emph{Define all work}. Decompose the effort using the \ac{wbs} and align it with the organizational structure. Control accounts are the natural management points where \acp{cam} are assigned, and each control account contains work packages and (if needed) planning packages.
  \item[Step 2:] \emph{Schedule the work}. Develop the integrated schedule (often visualized with Gantt charts) that sequences the \ac{wbs} elements, includes key milestones, and forms the backbone of the time-phased budget.
  \item[Step 3:] \emph{Budget the work}. Assign time-phased budgets to each work package, establish management reserve for in-scope known unknowns, and confirm that the sum of the control accounts equals the contract budget base.
\end{enumerate}
Changes to the \ac{pmb} must be formally controlled and documented. Reasons for changes include:
\begin{itemize}
  \item Contract changes
  \item Internal replanning
  \item Formal reprogramming
\end{itemize}

\subsection{EVM Reviews and Reports}
Formal reviews validate the integrity of \ac{evm} data products before and after system acceptance~\autocite{edo-3-6-1-evm-intro-2025}.
\subsubsection{Post-Acceptance Review}
\begin{itemize}
  \item Objective: ensure performance data remain accurate after system acceptance and identify any corrective actions required to reaffirm compliance.
  \item Timing: conducted as needed following system acceptance and prior to the next \ac{ibr}.
  \item Led by the Review Director (typically \ac{dcma}) with membership tailored to the purpose of the review.
  \item Culminates in a formal report prepared by the Review Director.
\end{itemize}
\subsubsection{Initial Compliance Evaluation}
\begin{itemize}
  \item Objective: validate that the Contractor's \ac{evm} system description matches actual practice and satisfies the \ac{evms} criteria.
  \item Timing: executed prior to the initial \ac{ibr} (and subsequently only if needed).
  \item Led by the \ac{dcma} Review Director with cross-functional support as required.
  \item Results documented in a report signed by the Review Director.
\end{itemize}
\subsubsection{Integrated Baseline Review}
\begin{itemize}
  \item Participants: the contractor, \ac{pm} and deputy, \ac{dcma}/\ac{dcaa}/\ac{supship} personnel, and relevant technical staff.
  \item Purpose: joint assessment to verify \ac{sow} coverage, logical schedule sequencing, resource realism, and risks captured in the \ac{pmb}.
  \item Timing: conducted no later than 180 calendar days after contract award, exercise of a significant option, or major modification in accordance with the \ac{dod} \ac{evm} Implementation Guide and service policy~\autocite{edo-3-6-1-evm-intro-2025,DFARS-EVMS-Deviation-2026}.
\end{itemize}
\subsubsection{Integrated Program Management Data and Analysis Report}
This contractually required report delivers cost, schedule, and technical status derived from the Contractor's \ac{evm}
system, enabling the \ac{pmo} to identify performance problem areas and emerging risks. When required, \ac{ipmdar} is
normally delivered monthly; final reports are due no later than 16 Government business days after the contractor
accounting period ends unless the contract tailors the deadline. The Rev.~D \ac{ipmdar} \ac{did} components
are~\autocite{IPMDAR-DID-RevD-2026,DoD-IPM-Policy-Guidance-2026,DoD-IPM-EVM-Definitions-2026}:
\begin{enumerate}
  \item Contract Performance Dataset.
  \item Schedule deliverables: the Schedule Performance Dataset and native \ac{ims}.
  \item Performance Narrative Report.
\end{enumerate}
A \ac{cfsr} complements the \ac{ipmdar} by providing funding and price information that reconciles to the cost-focused
earned value data~\autocite{edo-3-6-1-evm-intro-2025,DoD-IPM-EVM-Definitions-2026}.
\subsubsection{EVM Warning Signs}
When using \ac{evm}, the following are indications that a problem exists~\autocite{edo-3-6-2-evm-analysis-2025}.
\begin{itemize}
  \item Use of management reserves
  \item Significant revisions to the \ac{pmb}
  \item Zero variance
  \item Sudden change in trends
  \item Unreasonable \ac{tcpi}
\end{itemize}

\subsection{EVM Data Analysis}
\Fig{\ref{fig:evm_var}} shows the S-curve representation of the \ac{evm} variables; the gaps between the curves are the cost and schedule variances that drive performance assessments~\autocite{edo-3-6-2-evm-analysis-2025}.
\begin{figure}[H]
  \centering
  \includegraphics[width=0.7\linewidth]{EVM-chart.png}
  \caption[EVM Variances]{\ac{evm} chart showing the independent variables and variances. \srcCite{edo-3-6-2-evm-analysis-2025}}
  \label{fig:evm_var}
\end{figure}
\subsubsection{Data Analysis Steps}
\begin{enumerate}
  \item \textbf{Get current status}: capture the latest cost and schedule performance using \ac{cpi}, \ac{spi}, and narrative explanations.
  \item \textbf{Identify trends}: analyze indices and variance trajectories over time to highlight emerging risks and opportunities.
  \item \textbf{Predict completion}: develop independent \ac{eac} and \ac{vac} assessments to forecast final outcomes.
  \item \textbf{Determine management actions}: decide where to apply resources, contract changes, or risk mitigations to recover performance~\autocite{edo-3-6-2-evm-analysis-2025}.
\end{enumerate}
\subsubsection{Performance Measurement Techniques}
\begin{description}
  \item[Percent Complete:] Applies to long-duration tasks lacking interim milestones; progress is reported as the earned value fraction of total budget.
  \item[Weighted Milestones:] Uses milestone weights for long-duration tasks with clear waypoints, crediting earned value as milestones are completed.
  \item[Percent Start/Percent Finish:] Suitable for short-duration tasks (less than two reporting periods) with credit apportioned at start and finish.
  \item[0/100 Method:] Provides earned value only when the task is complete, offering a conservative measure for discrete, short tasks~\autocite{edo-3-6-2-evm-analysis-2025}.
\end{description}
\subsubsection{Equations}
Standard cost and schedule performance calculations include~\autocite{edo-3-6-2-evm-analysis-2025}.
\begin{align*}
  \text{CV}              & = \text{EV} - \text{AC}                                                                                 \\
  \text{SV}              & = \text{EV} - \text{PV}                                                                                 \\
  \text{CPI}             & = \frac{\text{EV}}{\text{AC}}                                                                           \\
  \text{SPI}             & = \frac{\text{EV}}{\text{PV}}                                                                           \\
  \text{CV} \%           & = \left(\frac{\text{CV}}{\text{EV}}\right)\times 100                                                    \\
  \text{SV} \%           & = \left(\frac{\text{SV}}{\text{PV}}\right)\times 100                                                    \\
  \text{\% Complete}     & = \left(\frac{\text{EV}}{\text{BAC}}\right)\times 100                                                   \\
  \text{\% Spent}        & = \left(\frac{\text{AC}}{\text{BAC}}\right)\times 100                                                   \\
  \text{EAC}             & = \frac{\text{BAC}}{\text{CPI}}                                                                         \\
  \text{EAC}             & = \text{AC} + \frac{\text{BAC} - \text{EV}}{\left(\frac{\text{EV}}{\text{AC}}\right)_{\text{3 months}}} \\
  \text{EAC}             & = \text{AC} + \frac{\text{BAC} - \text{EV}}{\text{CPI}\times\text{SPI}}                                 \\
  \text{EAC}             & = \text{AC} + \frac{\text{BAC} - \text{EV}}{0.8\,\text{CPI} + 0.2\,\text{SPI}}                          \\
  \text{VAC}             & = \text{BAC} - \text{EAC}                                                                               \\
  \text{TCPI}            & = \frac{\text{Work remaining}}{\text{Budget remaining}}                                                 \\
  \text{TCPI}_\text{EAC} & = \frac{\text{BAC} - \text{EV}}{\text{EAC} - \text{AC}}                                                 \\
  \text{TCPI}_\text{BAC} & = \frac{\text{BAC} - \text{EV}}{\text{BAC} - \text{AC}}                                                 
\end{align*}

\subsection{Quick Review}
\begin{itemize}
  \item What does \ac{evm} measure? \textbf{Answer:} Performance against the contract execution baseline by integrating
  scope, schedule, cost, and technical progress.
  \item What is the \ac{apb} versus \ac{pmb} trap? \textbf{Answer:} \ac{apb} is the government program baseline approved
  by the \ac{mdauthority}; \ac{pmb} is the contractor's time-phased baseline used for earned-value measurement.
  \item What do negative \ac{cv} and \ac{sv} mean? \textbf{Answer:} Negative \ac{cv} means cost overrun; negative
  \ac{sv} means schedule delay against planned value.
  \item Why does \ac{ibr} matter? \textbf{Answer:} It tests \ac{pmb} realism, scope coverage, schedule logic, resources,
  and risk within about 180 days after contract award, significant option, or major modification.
  \item What are warning signs? \textbf{Answer:} Management-reserve use, major \ac{pmb} revisions, zero variance, sudden
  trend shifts, and unreasonable \ac{tcpi}.
\end{itemize}

%====================
% End of file
%====================
\IfSubfilesClassLoaded{
  \RenewDocumentCommand{\entryname}{}{\textbf{\color{Modern} Acronym}}
  \RenewDocumentCommand{\descriptionname}{}{\textbf{\color{Modern} Definition}}
  \printnoidxglossary[
    type=\acronymtype,
    title=Acronyms,
    style=long-booktabs
  ]}{}
\end{document}
